% Copyright (C) 2013 - Michael Baudin
%
% This file must be used under the terms of the 
% Creative Commons Attribution-ShareAlike 3.0 Unported License :
% http://creativecommons.org/licenses/by-sa/3.0/

\chapter{The Halton sequence}

\section{Introduction}
    
The Hammersley sequence (1960) uses the Van Der Corput sequence
in order to create a low discrepancy point set \cite{Hammersley1960}.
However, the number of points in the Hammersley sequence must be know
in advance, which is sometimes an issue.
The Halton sequence (1964) uses the Van Der Corput sequence without
this limitation \cite{Halton1964}, and this is why it is used more often.

The Halton sequence in s dimensions uses s
different bases which must be relatively primes.
We usually use the s first prime numbers: 2, 3, 5, 7, 11, etc...

Denote by $p(i)$ the i-th prime number.
The dimension i of the Halton sequence is
the Van Der Corput in base $p(i)$.
Hence, the k-th element of the Halton sequence is
$$
  x_k=(\psi_{p_1}(k),\psi_{p_2}(k),...,\psi_{p_s}(k))
$$
for $k\geq 0$.

The following script produces 10 points of the Halton sequence
in 4 dimensions.
The prime number used here are: 2, 3, 5, 7.
    

\lstset{language=scilabscript}
\begin{lstlisting}
-->u=lowdisc_ldgen(9,4,"halton")
 u  =
    0.5       0.3333333    0.2     0.1428571  
    0.25      0.6666667    0.4     0.2857143  
    0.75      0.1111111    0.6     0.4285714  
    0.125     0.4444444    0.8     0.5714286  
    0.625     0.7777778    0.04    0.7142857  
    0.375     0.2222222    0.24    0.8571429  
    0.875     0.5555556    0.44    0.0204082  
    0.0625    0.8888889    0.64    0.1632653  
    0.5625    0.0370370    0.84    0.3061224  
\end{lstlisting}

The lowest components of the Halton sequence 
make use of smaller primes while higher components 
make use of the larger primes. 
Moreover, we know that the size of the cycle which 
allows the sequence to fill the space depends 
on the value of the corresponding prime. 
Hence, if the user of a Halton sequence orders the 
variables so that the most important variables come 
first, then the Halton sequence can potentially lead 
to better results. 

\section{Scrambled Halton Sequence}

In high dimensions, the Halton sequence poorly
fills the space, as can be seen in low dimensionnal projections.
In order to improve this, Braaten and Weller (1979) suggest to scramble
the digits in the Halton sequence, in order to break correlations\cite{braaten1979}.

The permuted radical inverse function is
$$
  \psi_b(k)=\sum_{i=0}^{r-1} \frac{\sigma(a_{i}(k))}{b^{i+1}}.
$$
where $\sigma$ is a given permutation.

For a given base b, the permutation does not change the values
of the points in the sequence:
only the ordering of the points changes.
However, since each dimension has a specific base b, the
multidimensionnal point set of a scrambled sequence is, in general,
different from the unscrambled point set.

The following function computes the scrambled Van Der Corput
radical inverse function.

\lstset{language=scilabscript}
\begin{lstlisting}
function r = scrambleVDC(index,b,sigma)
    current = index
    ib = 1.0 / b
    r = 0.0
    while (current>0)
        digit = modulo ( current , b )
        digit=sigma(digit+1)
        current = int ( current / b )
        r = r + digit * ib
        ib = ib / b
    end
endfunction
\end{lstlisting}

For example, the base-3 Van Der Corput sequence can be associated
with the permutation [0,2,1].
The following script computes the first ten points of
this scrambled sequence.

\lstset{language=scilabscript}
\begin{lstlisting}
sigma=[0,2,1];
b=3;
for index=1:10
    r = scrambleVDC(index,b,sigma);
    mprintf("index=%d, r=%f\n",index,r)
end
\end{lstlisting}

The previous script produces the following output.

\lstset{language=scilabscript}
\begin{lstlisting}
index=1, r=0.666667
index=2, r=0.333333
index=3, r=0.222222
index=4, r=0.888889
index=5, r=0.555556
index=6, r=0.111111
index=7, r=0.777778
index=8, r=0.444444
index=9, r=0.074074
index=10, r=0.740741
\end{lstlisting}

In their paper, Braaten and Weller provide the permutations
for dimensions up to 16.

\lstset{language=scilabscript}
\begin{lstlisting}
2 (0 1)
3 (0 2 1)
5 (0 3 1 4 2)
7 (0 4 2 6 1 5 3)
11 (0 5 8 2 10 3 6 1 9 7 4)
13 (0 6 10 2 8 4 12 1 9 5 11 3 7)
etc...
\end{lstlisting}

The issue here is that the limited size of the
table provided by Braaten and Weller limits the practical use of the
algorithm.
Tuffin (1996) provides improved algorithms to expand the table
(and to improve the permutations as well), but this work has
two limitations. 
First, the permutations are computed by taking all dimensions
into account, and not dimension-by-dimension as in
Braaten and Weller.
This has the advantage of improving the permutations, but has
the drawback of making the table larger.
Secondly, Tuffin published the algorithm, but not the actual permutations.

Kocis and Whiten (1979) provide a simple algorithm to generate
permutations for the Halton sequence \cite{Kocis1997}.
Their RR2 algorithm uses the Van Der Corput sequence in base 2 and
reverses the digits, removing the integers that are too large.
A more detailed description is presented in Lemieux \cite{Lemieux2009}.
This algorithm has no strong theoretical basis, but works
efficiently.

The following function computes the RR2 scrambling
for the i-th dimension of a point set with \scifun{s} dimensions,
using the bases in the matrix \scifun{b}.

\lstset{language=scilabscript}
\begin{lstlisting}
function sigma=RR2Scrambling(s, i, b)
    ns=ceil(log(b(s))/log(2))
    sigma=[];
    u=lowdisc_ldgen(2^ns,1,"halton")
    u=[0;u]
    for k=1:2^ns
        sigmak=u(k)*2^ns
        if (sigmak<b(i)) then
            sigma($+1)=sigmak;
        end
    end
endfunction
\end{lstlisting}

    
      The following script prints the RR2 scrambling
      used for a point set in 3 dimensions.
    

\lstset{language=scilabscript}
\begin{lstlisting}
s=3;
primemat=number_primes100();
b=primemat(1:s);
for i=1:s
    sigma=RR2Scrambling(s,i,b);
    sigmastr=strcat(string(sigma),",");
    mprintf("i=%d, [%s]\n",i,sigmastr);
end
\end{lstlisting}

    
      The previous script produces the following output.
    

\lstset{language=scilabscript}
\begin{lstlisting}
i=1, [0,1]
i=2, [0,2,1]
i=3, [0,4,2,1,3]
\end{lstlisting}

    
      This feature is provided by the "-scrambling" option of the
      \scifun{lowdisc\_configure} function.
